\documentclass[11pt,showkeys]{../../shared/essay}

\begin{document}

\title{\textit{Data for Whom:}\\
Responsible Data Collection in the Global South}

\author{Jaxen Dutta}
  \email[]{jaxen.dutta@uwaterloo.ca}
  \affiliation{David R. Cheriton School of Computer Science,\\
  University of Waterloo, Waterloo, Ontario, Canada}

\date{March 18, 2024}

\begin{abstract}
Technology companies from wealthy nations increasingly collect data from developing countries, presenting a double-edged opportunity: the potential for development impact alongside the risk of perpetuating extractive, neocolonial dynamics. This essay identifies six principles that should govern responsible data collection in these contexts: local context engagement and collaborative partnerships; informed consent with culturally adapted mechanisms; data minimization; local capacity building; robust data security adapted to infrastructure constraints; and fair benefit sharing with host communities. Underlying all six is a seventh, structural imperative: a conscious rejection of data colonialism, in which companies must operate as partners respecting local sovereignty rather than as external actors extracting a resource. Responsible data collection is not a one-time compliance exercise but an ongoing, iterative commitment to equitable engagement that serves the needs of the communities whose data is collected.
\end{abstract}

\keywords{data ethics, developing countries, informed consent, data colonialism, data minimization, capacity building, benefit sharing, global data governance, technology and development}

\maketitle

\section{Introduction}

In an increasingly data-driven world, technology companies from wealthy nations play a pivotal role in gathering information from developing countries. This process presents a double-edged sword: the potential for immense benefit intertwined with the risk of exploitation. Data collected in developing countries has been used to improve public health outcomes, inform infrastructure investment, and expand access to education and financial services~\cite{runde_innovative_2022}. Yet the same process, conducted without ethical grounding, can replicate colonial patterns of resource extraction in which value flows outward and communities are left without recourse or benefit~\cite{edwards_slave_2017}.

To ensure responsible data collection, companies must prioritize ethical considerations, respect for local communities, and a commitment to fostering equitable development. This essay explores six key principles and the structural imperative of avoiding data colonialism that technology companies should implement when collecting data in developing contexts.

\section{Local Context and Collaborative Engagement}

The foundation of responsible data collection lies in a deep understanding of the local context. Companies must move beyond a one-size-fits-all approach and actively engage with stakeholders on the ground. This necessitates collaboration with local governments to understand existing data collection frameworks and regulations, and open dialogue with community leaders and civil society organizations to grasp cultural sensitivities, access challenges, and potential concerns within the community~\cite{runde_innovative_2022}. Facilitating community meetings, focus groups, and participatory workshops can achieve this~\cite{grimmelikhuijsen_explaining_2022}.

Partnering with local universities and research institutions fosters knowledge exchange and leverages existing expertise. This collaborative approach ensures data collection methods are culturally appropriate and address the specific needs of the developing nation~\cite{goldberg_world_2020}. By building trust and transparency through these partnerships, companies can tailor data collection strategies to the unique challenges of the host country, ensuring the information gathered is relevant and impactful rather than merely convenient.

\section{Informed Consent and Privacy}

Respecting privacy and obtaining informed consent is paramount. Individuals in developing countries may have limited awareness of data privacy rights. To address this, companies must ensure data subjects understand the purpose and scope of data collection, including how their information will be used~\cite{ohm_broken_2009}. They should also be aware of potential risks associated with data collection, including data breaches or misuse, and must be empowered to freely choose participation, with clear opt-in and opt-out mechanisms readily available~\cite{ohm_broken_2009}.

Consent forms should be translated into local languages and presented in a clear and concise manner, avoiding technical jargon~\cite{jobin_global_2019}. Companies should consider alternative consent mechanisms beyond written forms to cater to populations with limited literacy or those who prefer oral communication. The \gls{gdpr}'s framework for meaningful, specific, and freely given consent offers a useful baseline standard even where it does not apply as a matter of law~\cite{eu_gdpr_2016}.

\section{Data Minimization}

Responsible data collection goes beyond simply obtaining consent. Companies must adhere to the principle of data minimization, collecting only the information necessary to achieve their stated objectives~\cite{aggarwal_general_2008}. Defining the purpose of data collection upfront and adhering to it ensures that information is used for legitimate reasons, such as improving public health services, education access, or infrastructure development~\cite{eu_gdpr_2016}. Unnecessary data should not be retained and should be securely deleted upon reaching its designated purpose. This principle protects individuals from downstream harms including surveillance, re-identification, and the misuse of information collected under one stated purpose for entirely different ends~\cite{ohm_broken_2009}.

\section{Local Capacity Building}

Data collection should not be a unidirectional extraction of resources. Technology companies have a responsibility to invest in building local capacity for sustainable development. This can be achieved through training programs offered in partnership with local institutions, equipping local professionals with data analysis, cybersecurity, and data privacy skills~\cite{hanna_digital_2017}. Empowering local professionals allows communities to manage their own data more effectively and participate actively in the data economy. Knowledge transfer and co-creation of data collection methods through collaboration with local universities and research institutions ensures that data collection addresses local priorities and fosters a sense of ownership within the community~\cite{runde_innovative_2022}.

\section{Data Security}

Data security is paramount. Companies must implement robust measures to safeguard data against breaches, unauthorized access, and cyber threats. This necessitates investing in secure servers, encryption technologies, and regular security audits. However, it is crucial to recognize limitations in some developing countries' infrastructure. Companies should adapt their methods and consider alternative data collection approaches when necessary~\cite{runde_innovative_2022}. Offline data collection methods, for example, may be required in regions with limited internet connectivity. The burden of security adaptation falls on the collecting company, not the host community: a company cannot discharge its security obligations simply by noting that local infrastructure is inadequate.

\section{Fair Compensation and Benefit Sharing}

Data is a valuable resource, and communities that contribute their information deserve a fair share of the benefits~\cite{runde_innovative_2022}. Companies should explore mechanisms for fair compensation and benefit sharing. This could involve investing in local development projects identified by the community itself, providing access to technology and digital skills training for local populations, or supporting local entrepreneurs and businesses that leverage data for social good~\cite{hanna_digital_2017}. This approach ensures that data collection contributes positively to the social and economic development of the host country rather than constituting pure extraction.

\section{Avoiding Data Colonialism}

Technology companies must be mindful of perpetuating neocolonial dynamics. Data colonialism, where external actors exploit local data resources for their own gain without adequate consent, compensation, or benefit to source communities, is a documented and growing concern~\cite{edwards_slave_2017}. The global landscape of \gls{ai} ethics guidelines is itself illustrative of this imbalance: Jobin et al.~\cite{jobin_global_2019} demonstrate that the jurisdictions producing these frameworks are overwhelmingly concentrated in the Global North, leaving the communities most affected by data extraction with the least voice in governing it.

Companies should collaborate as genuine partners, respecting local sovereignty and empowering communities to make informed decisions about their data~\cite{edwards_slave_2017}. This means working within existing legal frameworks, involving local stakeholders in data governance, and ensuring the data benefits the local population first and foremost~\cite{jobin_global_2019}. Ongoing surveillance of how data is subsequently used, and public accountability mechanisms, are essential components of this commitment.

\section{Long-Term Commitment and Continuous Improvement}

Responsible data collection is an ongoing process, not a one-time event. Companies must demonstrate a long-term commitment to collaborating with local stakeholders~\cite{runde_innovative_2022}. This commitment involves regular assessments to evaluate the impact of data collection initiatives and their alignment with community needs, establishing clear feedback channels that allow communities to voice their concerns and suggestions, and conducting periodic impact evaluations to assess the effectiveness of data collection in achieving its stated objectives. Utilizing feedback and evaluation results enables companies to continuously improve data collection practices and ensure they remain responsible and ethical over time.

\section{Conclusion}

By adhering to the principles of local context engagement, informed consent, data minimization, local capacity building, data security, and fair compensation, technology companies can foster responsible data collection practices in developing countries. Through collaborative partnerships and a commitment to long-term engagement, these companies have the potential to be powerful agents for positive change, leveraging data for sustainable development that prioritizes the needs and aspirations of the communities they serve. The alternative, data collection that extracts without reciprocating, is not merely an ethical failure: it is a continuation of a historical pattern of exploitation that the international development community has committed, at least in principle, to ending.

\bibliography{data-for-whom}
\end{document}
